\documentclass[a4paper,14pt]{extarticle}  % extarticle supports 14pt

% Encoding and language
\usepackage[utf8]{inputenc}
\usepackage[T1]{fontenc}
\usepackage[english]{babel}

% Fonts and headings
\usepackage{xcolor}
\usepackage{titlesec}
\titleformat{\section}
  {\bfseries\color{black}\LARGE}{\thesection}{1em}{}
\titleformat{\subsection}
  {\bfseries\color{black}\Large}{\thesubsection}{1em}{}
\usepackage{mathptmx}

% Page, tables, figures, maths
\usepackage{asymptote}
\usepackage{float}
\usepackage{geometry}
\usepackage{graphicx}
\usepackage{array}
\usepackage{booktabs}
\usepackage{caption}
\usepackage{amsmath,amssymb,amsthm}
\setcounter{secnumdepth}{0}
\geometry{
  top=2cm,
  bottom=2cm,
  left=3cm,
  right=2cm
}

% Paragraph indentation and line spacing
\setlength{\parindent}{1.25cm}
\renewcommand{\baselinestretch}{1.5}
\emergencystretch=2em % reduces overfull lines

% Boxed statements of theorems and lemmas (exactly as wide as the text)
\newsavebox{\thmboxbin}
\newenvironment{thmbox}
  {\par\medskip\noindent\begin{lrbox}{\thmboxbin}%
   \begin{minipage}{\dimexpr\linewidth-2\fboxsep-2\fboxrule\relax}}
  {\end{minipage}\end{lrbox}\fbox{\usebox{\thmboxbin}}\par\medskip}

% hyperref should be loaded last
\usepackage[colorlinks=true, urlcolor=blue, linkcolor=black, citecolor=black]{hyperref}

% Title
\title{Integer Solutions of the Diophantine Equation $x^3+y^3=x^2+kxy+y^2$}
\author{Maxud Eric}
\date{} % leave empty or put a date

\begin{document}
\maketitle
\section{Abstract}
Motivated by the Diophantine equation
\[
x^3+y^3=x^2+42xy+y^2,
\]
(BMO 2017/1)
\\
we study the generalized family
\[
x^3+y^3=x^2+kxy+y^2,
\qquad k\in\mathbb Z,
\]
over the integers. We prove that for every pair of positive integers \(n\) and \(l\), there exists an integer \(k>l\) such that the equation has more than \(n\) distinct integer solutions. On the other hand, we show that for every \(k\neq -1,2\), the number of integer solutions is finite and satisfies the explicit bound
\[
N_k \le 4\tau\!\left(|(k-2)^2(k+1)|\right),
\]
where \(\tau(m)\) denotes the number of positive divisors of \(m\). We also identify the exceptional cases \(k=-1\) and \(k=2\), for which the equation has infinitely many integer solutions.

\section{Introduction}
For small cases:
\begin{enumerate}
    \item if $k=0$, $$x^3+y^3=x^2+y^2$$ $$x^2-xy+y^2|x^2+y^2$$ Then $$x^2-xy+y^2|xy$$ So $x=y=0$ and $x=y=1$ are only solutions.
    \\
    \textbf{Answer:}
\[
(x,y) = (0,0), \ (1,1).
\]
    \item if $k=-1$, similarly we can get answers are all $x,y$ such that $x+y=1$ or $x^2-xy+y^2=0$ 
    \\
    \textbf{Answer:}
\[
(x,y) = (0,0), \ (x,1-x), \ (1-x,x).
\]
    \item if $k=1$, $$x^3+y^3=x^2+xy+y^2$$ Let
\[
s = x + y, \quad p = xy.
\]
Then
\[
x^3 + y^3 = (x+y)^3 - 3xy(x+y) = s^3 - 3ps,
\]
\[
x^2 + xy + y^2 = (x+y)^2 - xy = s^2 - p.
\]

Substituting, we get
\[
s^3 - 3ps = s^2 - p.
\]
Rearranging:
\[
s^3 - s^2 - 3ps + p = 0,
\]
\[
s^2(s - 1) - p(3s - 1) = 0.
\]
Thus,
\[
p(3s - 1) = s^2(s - 1).
\]

Observe that
\[
\gcd(3s - 1, s) = 1,
\]
since any common divisor divides
\[
(3s - 1) - 3s = -1.
\]

Also,
\[
\gcd(3s - 1, s - 1) \mid ((3s - 1) - 3(s - 1)) = 2.
\]

Hence from
\[
3s - 1 \mid s^2(s - 1)
\]
it follows that
\[
3s - 1 \mid 2.
\]

Therefore,
\[
3s - 1 \in \{-2, -1, 1, 2\}.
\]

Checking integer values:
\[
3s - 1 = -1 \Rightarrow s = 0, \quad
3s - 1 = 2 \Rightarrow s = 1.
\]

\textbf{Case 1:} $s = 0$  
\[
p(-1) = 0 \Rightarrow p = 0.
\]
Thus,
\[
x + y = 0, \quad xy = 0 \Rightarrow x = y = 0.
\]

\textbf{Case 2:} $s = 1$  
\[
p(2) = 0 \Rightarrow p = 0.
\]
Thus,
\[
x + y = 1, \quad xy = 0,
\]
so
\[
(x,y) = (1,0), \ (0,1).
\]

\textbf{Answer:}
\[
(x,y) = (0,0), \ (1,0), \ (0,1).
\]
\item
if \(k=2\), then the equation becomes
\[
x^3+y^3=x^2+2xy+y^2=(x+y)^2,
\]
and every pair with
\[
x+y=0
\]
is a solution, so there are infinitely many integer solutions.
\\
Besides, if $x+y \ne 0$ $$x^2+y^2=xy+x+y$$
Thus, if we double both sides: $$(x-y)^2+(x-1)^2+(y-1)^2=2$$
\\
\textbf{Answer:}
\[
(x,y) = (x,-x), \ (1,0), \ (0,1) \ (1,2) \ (2,1) \ (2,2).
\]
\end{enumerate}

\section{Theorems and Proofs}
For small $|k|$ the equation has only a few integer solutions (see the Introduction), yet Theorem~1 shows that the number of solutions is unbounded as $k$ varies.

\noindent\fbox{\begin{minipage}{\textwidth}
\paragraph{Theorem 1: (Almaty Math Battles 2026)} For every positive integers $n,l$, there is exist $k>l$, such that $$x^3+y^3=x^2+kxy+y^2$$ has more than $n$ solutions in $\mathbb{Z}$.
\end{minipage}}
\\
\\
\textbf{Proof.}
We look for solutions of the form
\[
y=mx
\]
for some positive integer \(m\). Substituting into
\[
x^3+y^3=x^2+kxy+y^2
\]
gives
\[
x^3+(mx)^3=x^2+kx(mx)+(mx)^2,
\]
hence
\[
x^3(1+m^3)=x^2(1+km+m^2).
\]
Therefore, if
\[
1+m^3 \mid 1+km+m^2,
\]
then
\[
x=\frac{1+km+m^2}{1+m^3},\qquad y=mx
\]
is an integer solution.

So for each chosen \(m\), it is enough to require
\[
1+km+m^2\equiv 0 \pmod{1+m^3}.
\]
A sufficient condition is
\[
k\equiv m(m-1)\pmod{1+m^3},
\]
because then
\[
1+km+m^2 \equiv 1+m^2(m-1)+m^2 = 1+m^3 \equiv 0 \pmod{1+m^3}.
\]

Now choose
\[
m_i=2^{2^i}\qquad (i=1,2,\dots,n+1),
\]
and let
\[
q_i=1+m_i^3=1+2^{3\cdot 2^i}.
\]

We claim that the integers \(q_1,q_2,\dots,q_{n+1}\) are pairwise coprime.  
Indeed, suppose \(i<j\) and some integer \(d>1\) divides both \(q_i\) and \(q_j\). Then
\[
2^{3\cdot 2^i}\equiv -1 \pmod d.
\]
Raising both sides to the power \(2^{j-i}\), we obtain
\[
2^{3\cdot 2^j}\equiv (-1)^{2^{j-i}}=1 \pmod d.
\]
On the other hand, since \(d\mid q_j\), we also have
\[
2^{3\cdot 2^j}\equiv -1 \pmod d.
\]
This is impossible. Hence
\[
\gcd(q_i,q_j)=1 \qquad \text{for all } i\ne j.
\]

Therefore, by the Chinese Remainder Theorem, there exists an integer \(k\) such that
\[
k\equiv m_i(m_i-1)\pmod{q_i}
\qquad (i=1,2,\dots,n+1).
\]
Moreover, all such integers form an arithmetic progression, so we may choose one with
\[
k>l.
\]

For each \(i\), define
\[
x_i=\frac{1+km_i+m_i^2}{1+m_i^3},\qquad y_i=m_i x_i.
\]
By construction, \(x_i\in\mathbb Z\), hence \((x_i,y_i)\in\mathbb Z^2\). Also,
\[
x_i^3+y_i^3=x_i^3(1+m_i^3)
= x_i^2(1+km_i+m_i^2)
= x_i^2+kx_i y_i+y_i^2,
\]
so \((x_i,y_i)\) is a solution of the equation.

These solutions are distinct, since
\[
\frac{y_i}{x_i}=m_i
\]
and the numbers \(m_1,m_2,\dots,m_{n+1}\) are pairwise distinct.

Thus the equation has at least \(n+1\) distinct integer solutions. In particular, it has more than \(n\) solutions in \(\mathbb Z\).

This proves the Theorem 1.

\noindent\fbox{\begin{minipage}{\textwidth}
\paragraph{Theorem 2:} Let
\[
N_k=\#\{(x,y)\in \mathbb Z^2:\ x^3+y^3=x^2+kxy+y^2\}.
\]
If \(k\neq -1,2\), then
\[
N_k\le 4\,\tau\!\left(\left|(k-2)^2(k+1)\right|\right),
\]
where \(\tau(m)\) denotes the number of positive divisors of \(m\).

Moreover, for \(k=-1\) and \(k=2\), the equation has infinitely many integer solutions.
\end{minipage}}
\\
\\
\textbf{Proof.}
Consider the equation
\[
x^3+y^3=x^2+kxy+y^2.
\]
Set
\[
s=x+y,\qquad p=xy.
\]
Then
\[
x^3+y^3=(x+y)^3-3xy(x+y)=s^3-3ps,
\]
and
\[
x^2+kxy+y^2=(x+y)^2+(k-2)xy=s^2+(k-2)p.
\]
Hence
\[
s^3-3ps=s^2+(k-2)p,
\]
so
\[
s^3-s^2=p(3s+k-2),
\]
that is,
\[
s^2(s-1)=p(3s+k-2).
\]
Let
\[
D=3s+k-2.
\]
Then
\[
D\mid s^2(s-1).
\]

We first note that if \(D=0\), then
\[
s^2(s-1)=0,
\]
so \(s=0\) or \(s=1\). Since \(D=3s+k-2\), this gives \(k=2\) or \(k=-1\). Therefore, for
\[
k\neq -1,2,
\]
we must have
\[
D\neq 0.
\]

We now prove that
\[
|D|\mid |(k-2)^2(k+1)|.
\]

Let \(q^a\parallel D\), where \(q\) is a prime. Since
\[
D\mid s^2(s-1)
\]
and
\[
\gcd(s,s-1)=1,
\]
the prime \(q\) divides exactly one of \(s\) and \(s-1\).

\medskip

\textit{Case 1: \(q\mid s-1\).}

Since \(q\nmid s\), from \(q^a\mid s^2(s-1)\) we get
\[
q^a\mid (s-1).
\]
But also \(q^a\mid D\), hence
\[
q^a\mid D-3(s-1)=k+1.
\]

\medskip

\textit{Case 2: \(q\mid s\).}

Write
\[
q^b\parallel s.
\]
Since \(q\nmid s-1\), from \(q^a\mid s^2(s-1)\) we obtain
\[
a\le 2b.
\]
Also, because \(q^a\mid D\) and \(q^b\mid s\), we have
\[
q^{\min(a,b)}\mid \gcd(D,s).
\]
But
\[
D-3s=k-2,
\]
so
\[
q^{\min(a,b)}\mid (k-2).
\]
If \(a\le b\), then \(q^a\mid (k-2)\). If \(a>b\), then \(q^b\mid (k-2)\), and since \(a\le 2b\), it follows that
\[
q^a\mid (k-2)^2.
\]
Thus in all cases,
\[
q^a\mid (k-2)^2.
\]

Therefore, every prime power divisor of \(D\) divides either \(k+1\) or \((k-2)^2\), and so
\[
|D|\mid |(k-2)^2(k+1)|.
\]

Now let
\[
M=(k-2)^2(k+1).
\]
Every integer solution \((x,y)\) gives rise to a divisor \(D\) of \(M\), and \(s=x+y\) is then uniquely determined by
\[
s=\frac{D-k+2}{3}.
\]
After that, \(p=xy\) is uniquely determined from
\[
p=\frac{s^2(s-1)}{D}.
\]

Finally, \(x\) and \(y\) are the roots of
\[
t^2-st+p=0.
\]
For fixed \(s\) and \(p\), this quadratic gives at most two ordered integer pairs \((x,y)\), namely \((x,y)\) and \((y,x)\).

Since \(M\neq 0\) for \(k\neq -1,2\), the number of integer divisors of \(M\) is
\[
2\tau(|M|).
\]
Hence
\[
N_k\le 2\cdot 2\tau(|M|)=4\tau(|M|).
\]
Therefore
\[
N_k\le 4\,\tau\!\left(\left|(k-2)^2(k+1)\right|\right).
\]

\bigskip

It remains to note the exceptional cases from introduction.
This completes the proof.

\noindent\fbox{\begin{minipage}{\textwidth}
\textbf{Lemma 1.} Let $p$ be a prime number. If $p$ can be represented in the form
\[
p = x^2 - xy + y^2
\]
for some integers $x,y$, then either $p=3$ or $p \equiv 1 \pmod{3}$.
\end{minipage}}
\\
\\
\textbf{Proof.}
\\
Consider the expression $x^2 - xy + y^2$ modulo $3$. Since every integer is congruent to $0,1,$ or $2$ modulo $3$, we analyze all possibilities for $x,y \pmod{3}$.

Note that for any integer $a$, we have $a^2 \equiv 0$ or $1 \pmod{3}$. A direct check shows that
\[
x^2 - xy + y^2 \equiv 0 \text{ or } 1 \pmod{3}
\]
for all integers $x,y$, and never $2 \pmod{3}$.

Now suppose that a prime $p$ satisfies
\[
p = x^2 - xy + y^2.
\]
Reducing modulo $3$, we obtain
\[
p \equiv 0 \text{ or } 1 \pmod{3}.
\]

If $p \equiv 0 \pmod{3}$, then $p=3$ since $p$ is prime. Otherwise, we must have
\[
p \equiv 1 \pmod{3}.
\]

\begin{thmbox}
\textbf{Theorem 3.} There exist infinitely many integers $k$ for which the equation
\[
x^3 + y^3 = x^2 + kxy + y^2
\]
has exactly four integer solutions (in particular, fewer than $5$). More precisely, if $p\neq 2$ is a prime with $p\equiv 2 \pmod 3$ and $k=p-1$, then the integer solutions are exactly
\[
(0,0),\quad (1,0),\quad (0,1),\quad \left(\tfrac{p+1}{2},\tfrac{p+1}{2}\right).
\]
\end{thmbox}

\begin{proof}
Let $p\ne 2$ be a prime with $p\equiv 2\pmod 3$ and let $k=p-1$. Note that $p$ is odd, so $t=\frac{p+1}{2}=\frac{k+2}{2}$ is an integer. The pairs $(0,0)$, $(1,0)$, $(0,1)$ are solutions for every $k$, and $(t,t)$ is a solution because $2t^3=(k+2)t^2$. These four pairs are distinct.

Conversely, let $(x,y)\ne(0,0)$ be an integer solution and let $d=\gcd(x,y)>0$. Write
\[
x = dx_1, \quad y = dy_1, \quad \gcd(x_1,y_1)=1.
\]
Substituting into the equation and dividing by $d^2$, we get
\[
d(x_1^3 + y_1^3) = x_1^2 + kx_1y_1 + y_1^2.
\]
Put $Q=x_1^2-x_1y_1+y_1^2>0$. Then $x_1^3+y_1^3=(x_1+y_1)Q$ and $x_1^2+kx_1y_1+y_1^2=Q+(k+1)x_1y_1$, so
\begin{equation}\label{eq:thm3}
d(x_1+y_1)Q = Q+(k+1)x_1y_1 .
\end{equation}
Hence $Q\mid (k+1)x_1y_1$. Now $\gcd(Q,x_1y_1)=1$: a prime dividing $x_1$ and $Q$ would divide $Q-x_1(x_1-y_1)=y_1^2$, hence $y_1$, contradicting $\gcd(x_1,y_1)=1$; the case of a prime dividing $y_1$ and $Q$ is symmetric. Therefore
\[
Q\mid k+1=p,
\]
so $Q\in\{1,p\}$. By Lemma~1, the prime $p\equiv 2\pmod 3$ is not of the form $x^2-xy+y^2$, hence $Q\neq p$ and so $Q=1$.

The equation $x_1^2-x_1y_1+y_1^2=1$ is equivalent to $(2x_1-y_1)^2+3y_1^2=4$, so $y_1\in\{-1,0,1\}$, and we find exactly six pairs:
\[
(x_1,y_1)\in\{(1,0),(-1,0),(0,1),(0,-1),(1,1),(-1,-1)\}.
\]
With $Q=1$ and $k+1=p$, equation~\eqref{eq:thm3} reads
\[
d(x_1+y_1)=1+p\,x_1y_1 .
\]
Recall $d>0$.
\begin{itemize}
    \item $(x_1,y_1)=(1,0)$: $d=1$, giving $(x,y)=(1,0)$.
    \item $(x_1,y_1)=(0,1)$: $d=1$, giving $(x,y)=(0,1)$.
    \item $(x_1,y_1)=(1,1)$: $2d=1+p$, giving $(x,y)=\left(\frac{p+1}{2},\frac{p+1}{2}\right)$.
    \item $(x_1,y_1)\in\{(-1,0),(0,-1),(-1,-1)\}$: the left-hand side is $-d$ or $-2d$ while the right-hand side is $1$ or $1+p$, both positive; this is impossible for $d>0$.
\end{itemize}
Together with $(0,0)$, these are exactly the four solutions listed.

Finally, there are infinitely many primes $p\equiv 2\pmod 3$ (a special case of Dirichlet's theorem; an elementary argument: if $p_1,\dots,p_r$ were all of them, then $3p_1\cdots p_r-1\equiv 2\pmod 3$ is coprime to $3$ and to every $p_i$, and not all of its prime factors can be $\equiv 1\pmod 3$, a contradiction). Hence we obtain infinitely many values $k=p-1$ with exactly four integer solutions.
\end{proof}


\begin{figure}[H]
\centering
\begin{asy}
// ------------------------------------------------------------
// Number of integer solutions of  x^3 + y^3 = x^2 + k*x*y + y^2
// as a function of k, for -50 <= k <= 50.
//
// Blue bars : exact number of integer (x,y) solutions for that k.
// Red bars  : k = -1 and k = 2, where the solution set is INFINITE
//             (a whole line of integer points satisfies the
//             equation). These are truncated at the top of the
//             frame and marked with "infinity", since their true
//             height is unbounded.
//
// Compile standalone with:
//     asy -f pdf solutions_graph.asy
//     asy -f png -render=4 solutions_graph.asy
// ------------------------------------------------------------

import graph;

size(15cm,11cm,IgnoreAspect);

int[] ks = {-50,-49,-48,-47,-46,-45,-44,-43,-42,-41,-40,-39,-38,-37,-36,-35,-34,-33,-32,-31,-30,-29,-28,-27,-26,-25,-24,-23,-22,-21,-20,-19,-18,-17,-16,-15,-14,-13,-12,-11,-10,-9,-8,-7,-6,-5,-4,-3,-2,-1,0,1,2,3,4,5,6,7,8,9,10,11,12,13,14,15,16,17,18,19,20,21,22,23,24,25,26,27,28,29,30,31,32,33,34,35,36,37,38,39,40,41,42,43,44,45,46,47,48,49,50};
real[] counts = {12.0,3.0,4.0,3.0,4.0,3.0,4.0,11.0,4.0,3.0,8.0,5.0,8.0,3.0,6.0,3.0,6.0,3.0,4.0,3.0,4.0,5.0,4.0,3.0,4.0,5.0,4.0,3.0,8.0,3.0,6.0,3.0,4.0,3.0,6.0,5.0,6.0,3.0,4.0,3.0,4.0,3.0,8.0,5.0,4.0,3.0,4.0,3.0,3.0,1.0,4.0,3.0,5.0,3.0,4.0,3.0,8.0,3.0,4.0,3.0,4.0,5.0,8.0,5.0,4.0,3.0,4.0,3.0,8.0,3.0,12.0,3.0,4.0,3.0,4.0,3.0,4.0,7.0,4.0,5.0,8.0,3.0,4.0,3.0,8.0,3.0,6.0,7.0,12.0,3.0,4.0,5.0,8.0,3.0,4.0,3.0,4.0,5.0,8.0,3.0,4.0};
bool[] isInf = {false,false,false,false,false,false,false,false,false,false,false,false,false,false,false,false,false,false,false,false,false,false,false,false,false,false,false,false,false,false,false,false,false,false,false,false,false,false,false,false,false,false,false,false,false,false,false,false,false,true,false,false,true,false,false,false,false,false,false,false,false,false,false,false,false,false,false,false,false,false,false,false,false,false,false,false,false,false,false,false,false,false,false,false,false,false,false,false,false,false,false,false,false,false,false,false,false,false,false,false,false};

real barw     = 0.8;
real ymax     = 14;    // frame height; true finite counts top out at 12
real infTop   = ymax;

pen normalFill = rgb(0.20,0.45,0.75);
pen infFill    = rgb(0.85,0.20,0.20);
pen thinBlack  = black+linewidth(0.3pt);

for (int i = 0; i < ks.length; ++i) {
  real xpos   = ks[i];
  real h      = isInf[i] ? infTop : counts[i];
  pen fillCol = isInf[i] ? infFill : normalFill;

  filldraw(box((xpos-barw/2,0),(xpos+barw/2,h)), fillCol, thinBlack);

  if (isInf[i]) {
    label("$\infty$", (xpos, ymax+0.6), p=infFill+fontsize(11pt));
  }
}

// reference line: (0,0),(1,0),(0,1) solve the equation for every k
draw((-52,3)--(52,3), dashed+gray(0.5));
label("baseline = 3 ",
      (-15,12), p=gray(0.4)+fontsize(12pt));

xaxis("$k$", BottomTop, LeftTicks(Step=10, step=5));
yaxis("number of integer solutions", LeftRight, RightTicks(Step=2, step=1));

label("Integer solutions of $x^3+y^3 = x^2+kxy+y^2$ for $-50 \le k \le 50$",
      (0, ymax+1.8), p=fontsize(12pt));

\end{asy}
\caption{The number $N_k$ of integer solutions for $-50\le k\le 50$. The red bars mark $k=-1$ and $k=2$, where $N_k=\infty$; the dashed line is the baseline $N_k\ge 3$ given by the trivial solutions.}
\label{fig:counts}
\end{figure}

Figure~\ref{fig:counts} shows $N_k$ for $-50\le k\le 50$. For $k\ne -1,2$ in this range we have $3\le N_k\le 12$, and $N_k\in\{3,4\}$ for $66$ of the $99$ values of $k$. The trivial estimate $\tau(m)\le 2\sqrt m$ turns Theorem~2 into $N_k=O\!\left(|k|^{3/2}\right)$, which is much weaker than what the data suggest. A finer estimate of the divisor function gives a linear bound; the proof uses two small lemmas.

\input{Lemma_2_1}

\input{Lemma_2_2}

\begin{thmbox}
\textbf{Theorem 4.} There exists an integer constant $c$ such that for every integer $k$ with $|k|\ge 3$ the equation
\[
x^3+y^3=x^2+kxy+y^2
\]
has at most $c\,|k|$ integer solutions, i.e.\ $N_k\le c\,|k|$.
\end{thmbox}

\begin{proof}
Let $|k|\ge 3$; in particular $k\neq -1,2$. Put $M=\left|(k-2)^2(k+1)\right|\ge 1$ and write $M=\prod_i p_i^{a_i}$ with distinct primes $p_i$. Then $\tau(M)=\prod_i(a_i+1)$ and $M^{1/3}=\prod_i p_i^{a_i/3}$, so
\[
\frac{\tau(M)}{M^{1/3}}=\prod_i f(p_i,a_i).
\]
By Lemma~2.1 every factor with $p_i>27$ is at most $1$, and by Lemma~2.2 every factor with $p_i<27$ is at most $C_{p_i}$. Since all $C_p\ge 1$, we get
\[
\tau(M)\le C\,M^{1/3},\qquad C:=\prod_{p<27}C_p .
\]
For $|k|\ge 3$ we have $|k-2|\le \frac53|k|$ and $|k+1|\le\frac43|k|$, hence $M\le \frac{100}{27}|k|^3$ and $M^{1/3}\le \left(\frac{100}{27}\right)^{1/3}|k|$. By Theorem~2,
\[
N_k\le 4\tau(M)\le 4C\left(\tfrac{100}{27}\right)^{1/3}|k| .
\]
Thus the theorem holds with any integer $c\ge 4C\left(\frac{100}{27}\right)^{1/3}$.
\end{proof}

\noindent\textit{Remarks.} (1) By the explicit values above, $C\approx 3.53$, so $c=22$ works.
(2) Replacing the exponent $\frac13$ by an arbitrary $\varepsilon>0$ in the definition of $f$, the same argument gives $N_k=O_\varepsilon\!\left(|k|^{\varepsilon}\right)$ for every $\varepsilon>0$. So Theorem~4 is far from sharp as an upper bound (for example we can use Wigert's works to obtain better upper bound); by Theorem~1, however, $N_k$ is unbounded. 


\section{Conclusion}
For equations of this form, it seems difficult to obtain a complete general description of all integer solutions. Nevertheless, the results proved above show that for every \(k\neq -1,2\), the equation has only finitely many integer solutions, while for the exceptional values \(k=-1\) and \(k=2\), the set of integer solutions is infinite. Moreover, the number of solutions is unbounded as \(k\) varies (Theorem~1), but it grows at most linearly in \(|k|\) (Theorem~4), and Theorem~3 shows that for infinitely many integers \(k\) the equation has exactly four integer solutions.

\newpage
\begin{thebibliography}{9}
\bibitem{siu} \href{https://bibnum.publimath.fr/ACF/ACF19075.pdf}{Equations in China: Two Millennia of Innovation, Transmission and Re-transmission}, Man-Keung.
\bibitem{bmo} \href{https://artofproblemsolving.com/community/c6h1441692p8209533}{Balkan Mathematical Olympiad 2017, Problem 1} (Art of Problem Solving thread).
\bibitem{prime} \href{https://www.awesomemath.org/wp-pdf-files/math-reflections/mr-2025-03/mr_3_2025_on_prime_numbers_of_the_form.pdf}{On Prime Numbers of the Form $a^2 + ab + b^2$ Where $a \in \mathbb{N}^{+}$, $b \in \mathbb{N}^{+}$}, Arno Mikaelyan, Nairi Sedrakyan, Hayk Sedrakyan, Mathematical Reflections, 2025, no.~3.
\bibitem{tau} \href{https://www.renyi.hu/~gharcos/wigert4.pdf}{ON THE HYBRID MOMENTS OF THE GENERALIZED DIVISOR FUNCTION}, Gergely Harcos
\end{thebibliography}
\end{document}
